% =====================================================================
% HMM: численный пример Forward
% =====================================================================
\begin{frame}{HMM Forward: численный пример}

\textbf{Партитура:} 3 ноты $P = (60, 62, 64)$.

\textbf{Параметры модели:}
\begin{itemize}
\item Начальное: $\pi = (1, 0, 0)$ (стартуем с первой ноты).
\item Переходы: $a_{j, j} = 0.3$ (\textit{stay}),
      $a_{j, j+1} = 0.6$ (\textit{advance}),
      $a_{j, j+2} = 0.1$ (\textit{skip}).
\item Эмиссия: гауссиана $b_i(o) = \frac{1}{\sigma\sqrt{2\pi}}
      \exp\!\bigl(-\tfrac12 (\tfrac{o - p_i}{\sigma})^2\bigr)$,
      $\sigma = 2$.
\end{itemize}

\textbf{Поток:} $O = (60, 62)$.

\bigskip

\textbf{Шаг $t=1$, $o_1 = 60$.} Считаем эмиссии:
\[
b_1(60) \!=\! \tfrac{1}{2\sqrt{2\pi}} \!\approx\! 0.199,\quad
b_2(60) \!=\! 0.121,\quad
b_3(60) \!=\! 0.027.
\]
По формуле $\alpha_1(i) = \pi_i\, b_i(o_1)$ получаем
$\alpha_1 = (0.199,\ 0,\ 0)$. Нормировка: $\alpha_1 = \rlhi{(1, 0, 0)}$.

\textbf{Предсказание:} $\hat\imath_1 = 1$, $\mathrm{conf}_1 = 1.0$.

\end{frame}

\begin{frame}{HMM Forward: численный пример (продолжение)}

\textbf{Шаг $t=2$, $o_2 = 62$.}

\textbf{Шаг A.} Подсчёт prior $\sum_j \alpha_1(j)\, a_{j, i}$:
\[
\text{prior}_1 = 1 \cdot 0.3 = 0.3, \quad
\text{prior}_2 = 1 \cdot 0.6 = 0.6, \quad
\text{prior}_3 = 1 \cdot 0.1 = 0.1.
\]

\textbf{Шаг B.} Эмиссии $b_i(62)$:
\[
b_1(62) \!\approx\! 0.121, \quad
b_2(62) \!=\! 0.199, \quad
b_3(62) \!\approx\! 0.121.
\]

\textbf{Шаг C.} Умножение и нормировка:
\[
\widetilde{\alpha}_2 = (0.0363,\ 0.1194,\ 0.0121),
\quad
\alpha_2 = \rlhi{(0.217,\ 0.711,\ 0.072)}.
\]

\bigskip

\textbf{Предсказание:} $\hat\imath_2 = \argmax \alpha_2 = \rlhi{2}$,\quad
$\mathrm{conf}_2 = 0.711$.

\bigskip

\textbf{Интерпретация:} модель \emph{уверена} что мы во второй ноте
(71\,\%), небольшая масса остаётся на первой (мы могли ошибиться) и
третьей (могли skip-нуть).

\textbf{Этот вектор $\alpha_t$ --- основной сигнал для RL-агента.}

\end{frame}
